\section{Introduction}
\label{sec:introduction}

Parameter-efficient fine-tuning (PEFT) is a practical mechanism for adapting
large language models without updating the pretrained backbone
\citep{hu2022lora,houlsby2019adapter}. Its continual use, however, creates a
resource question that is absent from one-shot fine-tuning: how should a fixed
low-dimensional update space accommodate a potentially unbounded stream of new
data? Reusing one adapter can overwrite previously useful behavior, whereas
storing one adapter per stage makes memory and deployment complexity grow with
the stream.

Much of continual-PEFT evaluation presents semantically distinct tasks one at a
time and either stores task-specific modules or assumes that a task identity is
available at inference. This abstraction is useful for diagnosis but is not the
only deployment pattern. A production model may instead receive successive
batches of documents, demonstrations, or feedback drawn from an evolving
mixture of sources. The prediction interface and loss remain unchanged, source
boundaries can be unknown, and only one adapter is deployed. We call this the
\emph{task-identity-free, clocked temporal-mixture} setting. A deployment clock
or fixed minibatch schedule determines when an update is attempted, but no
semantic task or source identity is supplied. This setting includes gradual
domain drift, abrupt and recurrent mixtures, expanding support, and
chronological knowledge updates; it is not the stronger boundary-free setting
in which the learner must also decide when a distribution change occurred.

Recent analysis of scale sharpens why this setting is hard. \citet{huang2026larger}
show that larger models retain rare and complex tasks through \emph{reduced
interference}: once enough capacity is devoted to common tasks, the common-task
gradients weaken and stop overwriting the slowly accumulating features of rare
tasks. Their mechanism identifies retention with capacity, and their remedy is
scale---train a wider model. Continual PEFT is precisely the regime in which
that remedy is unavailable: the backbone is frozen, one adapter is deployed, and
the total rank budget is fixed and does not grow with the stream. The binding
constraint here is not parameter count---full fine-tuning also keeps its
parameter count fixed---but the \emph{effective update dimension relative to the
diversity of tasks the stream demands}. Low-rank adaptation makes this dimension
an explicit, tunable budget $R$; full fine-tuning is the slack limit in which
that budget is large enough that the capacity term stops binding and forgetting
reduces to conflict and online error. Our effective-capacity construction
(Proposition~\ref{prop:effective-capacity}) lets us compare the two empirically
at matched effective rank; we treat this as a comparator, not as a theorem
unifying the two training regimes.
Freezing the backbone is what keeps this analysis clean: historical curvature
then accumulates in stable ambient coordinates, whereas full fine-tuning rotates
the representation itself and moves those coordinates. This motivates
our central question. We do \emph{not} claim that allocation can manufacture the
capacity that scale buys: a fixed budget carries an unavoidable spectral-tail
capacity floor (Theorem~\ref{thm:spectral-capacity}), and no allocation rule can
remove it. The question is instead where the remaining, avoidable loss lives:
\emph{under a fixed budget, can explicit risk-reweighted reallocation reduce the
avoidable online allocation error and shift the retention--plasticity trade-off
so that the worst-served source---the one scale would have protected---is
retained, at a controlled cost to frequent-task retention?} We give controlled evidence that
under a fixed budget the scale-driven protection of the worst-served source does
not appear on its own---every allocator that spends the budget \emph{implicitly}
charges the aggregate capacity floor to the worst-served source, so their
worst-group retention clusters at a common low value---and that an
\emph{explicit} risk-reweighted re-solve \emph{redistributes} the same fixed
budget across sources, raising worst-group retention at a matched
frequency-weighted mean rather than lowering the aggregate floor or reducing
forgetting uniformly.

Answering this requires treating continual forgetting, in part, as a
capacity-allocation problem in a shared low-dimensional update space. The word
``partly'' is essential. Rank cannot explain incompatible target behavior: a
single unconditional model may be unable to satisfy two contradictory
distributions even with unlimited adapter rank. Nor does sequential training
always produce monotone forgetting, because later updates can cancel earlier
ones or produce positive backward transfer. A defensible capacity claim must
therefore separate three effects: irreducible conflict among distributions,
approximation error of the fixed-rank model class, and error from allocating and
optimizing that class online.

We develop this separation under a local second-order model. A Taylor expansion
shows precisely when the linear term in old loss disappears: the earlier
checkpoint must be stationary for the same historical risk in the feasible
update directions. Under a quadratic old loss, forgetting is the curvature
energy of the sum of subsequent updates. This yields accumulation in
expectation under a stated non-cancellation or stochastic-innovation condition,
but also exposes simple cancellation counterexamples. At the representation
level, an exact telescoping identity decomposes the historical retention gap
into conflict, fixed-rank capacity, and online error. For a Kronecker-factored
quadratic loss, the capacity term is exactly a curvature-whitened spectral tail.
Finally, as positive-semidefinite historical curvature is accumulated, its
nullspace---the exact second-order zero-interference space---can only contract.

The analysis suggests a solution other than continually increasing rank.
\method{} (Pareto-Safe Rank Reallocation) keeps a single fixed-rank-$R$ adapter
and, at each window, re-solves a risk-reweighted reduced-rank regression from a
bounded replay buffer. The distinction our evidence turns on is
\emph{implicit} versus \emph{explicit} allocation. Dense updating, replay, and
gradient-projection defenses allocate the $R$ retained directions implicitly:
the subspace they keep is whatever the accumulated gradient mass spans, which the
frequent sources dominate, so the worst-served source is left in the spectral
tail. \method{} instead spends the fixed memory budget on \emph{coverage}, keeping a
small dual buffer whose diverse part retains rare input directions that a single
uniform reservoir under-samples, and then runs a minimax multiplicative-weights
loop that moves the weight onto the worst-served buffered windows and re-solves
the closed-form weighted rank-$R$ problem, driving the $R$-th direction toward
that source; it keeps the iterate with the smallest buffered worst-window loss,
so the update is never worse than the uniform solve \emph{on that buffered
worst window} by construction, with no rollback or trust region required; the
name's ``safe'' denotes exactly this buffer-level property, not a held-out or
Pareto guarantee. The coverage buffer is an empirically
motivated remedy for rare-source under-coverage, not a consequence of the
capacity theorem; the weighted re-solve it feeds is the theorem-tied step. The
learner reads only the window data---no source label, task identifier, generator
oracle, or router---and the deployed rank never grows.

Our contributions are:
\begin{itemize}
  \item We formulate task-identity-free, clocked temporal-mixture continual
  PEFT with one deployed adapter, a fixed total rank budget, a smaller
  per-window active-rank budget, and bounded historical memory.
  \item We give a conditional theorem chain that (i) expands signed forgetting,
  (ii) characterizes sequential accumulation and its cancellation failure mode,
  (iii) exactly separates conflict, capacity, and online error, and (iv) proves
  contraction of the locally safe update space.
  \item We turn the exact rank-capacity characterization into an online method,
  \method{}: a bounded-buffer, minimax risk-reweighted reduced-rank regression
  whose per-window solve is the curvature-whitened truncated SVD of
  Theorem~\ref{thm:spectral-capacity}, and which is safe against the uniform
  solve on the buffered worst window by construction (a buffer-level property,
  not a held-out guarantee). The same spectral result gives a precise but limited
  connection to curvature-aware model merging.
  \item We recast the reduced-interference mechanism of \citet{huang2026larger}
  as an allocation question under a fixed budget, and give controlled evidence,
  under an internally pre-specified, hash-recorded held-out test on a crossed
  design of held-out geometries and stream draws, that implicit allocators (dense
  updating, reservoir/GSS/MIR replay, gradient projection, and the uniform-weight
  solve) all charge the aggregate spectral-tail floor to the worst-served source,
  clustering at a common worst-group retention, whereas coverage-aware explicit
  minimax reallocation \emph{redistributes} the fixed budget to raise worst-group
  retention at a matched frequency-weighted mean---a paired-average gain with
  cluster-robust confidence intervals excluding zero, not uniform dominance and
  not a lowering of the aggregate floor.
  \item We specify controlled mixture-drift and chronological experiments that
  test whether spectral tails and safe-space contraction predict forgetting,
  with total rank, active rank, replay, and optimizer-state budgets held equal
  across arms and compute reported per arm rather than matched---the method
  deliberately trades extra compute (repeated closed-form re-solves) for tail
  retention, which we disclose in a resource ledger rather than hide.
\end{itemize}
